% Terjemahan Bahasa Indonesia dari chapter2.tex baris 1133--1244.
% Sumber: Methods of Algebra, Volume 2, commit
% 9a5803ff77dd3257484cb177f851a73770a59dd3 (CC BY 4.0).
% stable-unit-id: o014.aljabr2.chapter2.simplicity-and-semisimplicity

% segment-id: o014.li2.chapter2.simplicity-and-semisimplicity.h001
\section{Objek Sederhana dan Semisederhana}\label{sec:semisimple}
\hypertarget{o014.aljabr2.chapter2.simplicity-and-semisimplicity}{ }

% segment-id: o014.li2.chapter2.simplicity-and-semisimplicity.p001
Pada bagian ini, tetapkan suatu kategori abelian $\mathcal{A}$.

% segment-id: o014.li2.chapter2.simplicity-and-semisimplicity.c001
\begin{konvensi}\label{con:direct-sum}
    \index{jumlah langsung (direct sum)}
    Secara konvensi, koproduk $\coprod_{i \in I} X_i$ dari suatu keluarga
objek $(X_i)_{i \in I}$ dalam kategori abelian (dengan asumsi koproduk itu ada)
ditulis dalam bentuk $\bigoplus_{i \in I} X_i$ dan disebut
\emph{jumlah langsung}. Jika $I$ berhingga, semua konvensi dari
Definisi~\ref{def:additive-category} berlaku.
\end{konvensi}

% segment-id: o014.li2.chapter2.simplicity-and-semisimplicity.d001
\begin{definisi}\index{objek sederhana (simple object)}
    Jika objek $X$ dalam $\mathcal{A}$ tidak nol dan
$\mathrm{Sub}_X = \left\{0, X \right\}$, maka $X$ disebut
\emph{objek sederhana}.
\end{definisi}

% segment-id: o014.li2.chapter2.simplicity-and-semisimplicity.p002
Objek $X$ sederhana berarti bahwa untuk setiap barisan eksak pendek
$0 \to X' \to X \to X'' \to 0$, tepat satu di antara $X'=0$ dan
$X'\rightiso X$ berlaku; secara ekuivalen, tepat satu di antara
$X\rightiso X''$ dan $X''=0$ berlaku. Karena itu, kesederhanaan $X$ dalam
$\mathcal{A}$ ekuivalen dengan kesederhanaannya dalam $\mathcal{A}^{\opp}$.

% segment-id: o014.li2.chapter2.simplicity-and-semisimplicity.p003
Perhatikan bahwa $X\neq0$ jika dan hanya jika $\identity_X\neq0$, dan ini
juga ekuivalen dengan pernyataan bahwa $\End(X)$ merupakan gelanggang tak nol.

% segment-id: o014.li2.chapter2.simplicity-and-semisimplicity.lm001
\begin{lema}[Lema Schur]
    Ambil objek $X,Y$. Jika $X$ (atau $Y$) merupakan objek sederhana, maka
semua morfisme tak nol dalam $\Hom(X,Y)$ merupakan monomorfisme (atau
epimorfisme). Sebagai akibatnya, jika $X$ merupakan objek sederhana, maka
$\End(X)$ merupakan gelanggang pembagian.
\end{lema}
% segment-id: o014.li2.chapter2.simplicity-and-semisimplicity.pf001
\begin{bukti}
    Misalkan $X$ sederhana dan $f:X\to Y$ tak nol. Maka haruslah
$\Ker(f)=0$. Kasus ketika $Y$ sederhana ditangani dengan dualitas. Morfisme
yang sekaligus merupakan monomorfisme dan epimorfisme adalah isomorfisme
(Proposisi~\ref{prop:strict-isomorphism}); maka, dengan mengambil $X=Y$,
diperoleh bahwa $\End(X)$ merupakan gelanggang pembagian.
\end{bukti}

% segment-id: o014.li2.chapter2.simplicity-and-semisimplicity.p004
Telah diketahui bahwa $\mathrm{Sub}_X$ merupakan kisi modular terbatas
(Teorema~\ref{prop:subobject-modularity}). Lema~\ref{prop:lattice-composition-series}
menunjukkan bahwa $X$ berpanjang hingga jika dan hanya jika himpunan terurut
parsial $[0,X]$ mempunyai deret komposisi; deret tersebut juga disebut
deret komposisi dari $X$. Untuk objek $X$ berpanjang hingga, \emph{panjang}-nya
didefinisikan sebagai \index{panjang (length)}
\index[sym1]{lX@$\ell(X)$}
% segment-id: o014.li2.chapter2.simplicity-and-semisimplicity.q001
\[ \ell(X) := \text{panjang dari }\mathrm{Sub}_X\;\in \Z_{\geq 0}; \]
lihat Definisi~\ref{def:lattice-length}. Perhatikan bahwa
$\ell(X)=0\iff X=0$. Berdasarkan konvensi, $Y\subsetneq Z$ berarti
$Y\subset Z$ dan $Y\neq Z$.

% segment-id: o014.li2.chapter2.simplicity-and-semisimplicity.dt001
\begin{definisiteorema}[Teorema Jordan--Hölder]\label{def:JH}
    \index{Teorema Jordan-Holder@Teorema Jordan--Hölder (Jordan--Hölder Theorem)}
\index{faktor komposisi (composition factor)}
\index[sym1]{JH@$\mathrm{JH}(X)$}
\sloppy
Misalkan $X$ objek berpanjang hingga. Ambil deret komposisi
$X=X_0\supsetneq\cdots\supsetneq X_r=0$. Objek $X_i/X_{i+1}$, yang dipandang
hingga isomorfisme, disebut \emph{faktor komposisi}. Faktor-faktor ini sederhana.
Multihimpunannya (dengan multiplisitas) dilambangkan dengan $\mathrm{JH}(X)$
dan tidak bergantung pada deret yang dipilih.
\fussy
\end{definisiteorema}
% segment-id: o014.li2.chapter2.simplicity-and-semisimplicity.pf002
\begin{bukti}
    Faktor komposisi pasti sederhana; jika tidak, $X_0\supsetneq\cdots\supsetneq
X_r$ mempunyai penghalusan sejati. Misalkan
$X=X'_0\supsetneq\cdots\supsetneq X'_s$ merupakan deret komposisi lain.
Teorema~\ref{prop:lattice-JH} menyatakan bahwa deret ini ekuivalen dengan
$X_0\supsetneq\cdots\supsetneq X_r$; khususnya, terdapat bijeksi
$i\leftrightarrow j$ antara himpunan indeksnya sedemikian sehingga interval
$\left[X_{i+1},X_i\right]$ dan $\left[X'_{j+1},X'_j\right]$ dapat dihubungkan
oleh isomorfisme standar dalam kisi modular,
$[a\wedge b,a]\rightiso[b,a\vee b]$ (lihat
Proposisi~\ref{prop:lattice-diamond-isom}).

    Sampai di sini semuanya dinyatakan dalam bahasa teori kisi. Namun,
Teorema~\ref{prop:Abel-cat-isom-thm} (iii) menunjukkan bahwa isomorfisme
interval semacam ini bersesuaian dengan isomorfisme objek kuosien dalam
$\mathcal{A}$,
% segment-id: o014.li2.chapter2.simplicity-and-semisimplicity.q002
\[ X_i/X_{i+1} \simeq X'_j/X'_{j+1}. \]
Dengan memvariasikan pasangan $i\leftrightarrow j$, terlihat bahwa, hingga isomorfisme
dan pengurutan ulang, faktor-faktor komposisi tidak bergantung pada pilihan
deret komposisi.
\end{bukti}

% segment-id: o014.li2.chapter2.simplicity-and-semisimplicity.p005
Perhatikan bahwa banyaknya elemen dalam $\mathrm{JH}(X)$ (dihitung dengan
multiplisitas) tepat sama dengan $\ell(X)$. Multihimpunan juga memiliki operasi
gabungan, yang berarti penjumlahan multiplisitas; operasi ini tetap
dilambangkan dengan $\cup$.

% segment-id: o014.li2.chapter2.simplicity-and-semisimplicity.lm002
\begin{lema}
    Diberikan barisan eksak pendek $0\to X'\to X\to X''\to0$. Objek $X$
berpanjang hingga jika dan hanya jika $X'$ dan $X''$ juga demikian; dalam hal
ini, $\mathrm{JH}(X)=\mathrm{JH}(X')\cup\mathrm{JH}(X'')$, sehingga
$\ell(X)=\ell(X')+\ell(X'')$.
\end{lema}
% segment-id: o014.li2.chapter2.simplicity-and-semisimplicity.pf003
\begin{bukti}
    Teorema~\ref{prop:Abel-cat-isom-thm} (ii) mengidentifikasi himpunan terurut
parsial $\mathrm{Sub}_{X'}$ dan $\mathrm{Sub}_{X''}$ masing-masing dengan
interval $[0,X']$ dan $[X',X]$ dalam $\mathrm{Sub}_X$. Jadi, jika $X$
berpanjang hingga, maka $X'$ dan $X''$ juga berpanjang hingga. Sebaliknya,
misalkan $X'$ dan $X''$ berpanjang hingga, dan ambil deret komposisi
% segment-id: o014.li2.chapter2.simplicity-and-semisimplicity.q003
\[ X'=X'_0\supsetneq\cdots\supsetneq X'_r=0,\quad
X''=X''_0\supsetneq\cdots\supsetneq X''_s=0. \]
Menurut Teorema~\ref{prop:Abel-cat-isom-thm} (ii), angkat setiap $X''_i$
menjadi subobjek $Y_i$ dari $X$ dengan $Y_i\supset X'$, khususnya $Y_0=X$
dan $Y_s=X'$; maka
% segment-id: o014.li2.chapter2.simplicity-and-semisimplicity.q004
\[ X=Y_0\supsetneq\cdots\supsetneq Y_s\supsetneq X'_1\supsetneq\cdots
\supsetneq X'_r=0 \]
merupakan deret komposisi, yang subkuosiennya membentuk
$\mathrm{JH}(X')\cup\mathrm{JH}(X'')$.
\end{bukti}

% segment-id: o014.li2.chapter2.simplicity-and-semisimplicity.d002
\begin{definisi}\label{def:semisimple}
    \index{objek!semisederhana (semisimple)}
    \index{objek!terbelah (split)}
    \index{kategori abelian!semisederhana (semisimple)}
    \index{kategori abelian!terbelah (split)}
    Objek $X$ disebut
    % segment-id: o014.li2.chapter2.simplicity-and-semisimplicity.l001
    \begin{itemize}
        % segment-id: o014.li2.chapter2.simplicity-and-semisimplicity.i001
        \item \emph{semisederhana} jika
        $X=\bigoplus_{i\in I}X_i$ untuk suatu keluarga objek sederhana
        $(X_i)_{i\in I}$;
        % segment-id: o014.li2.chapter2.simplicity-and-semisimplicity.i002
        \item \emph{terbelah} jika setiap barisan eksak pendek
        $0\to X'\to X\to X''\to0$ terbelah.
    \end{itemize}
    Jika semua objek dalam $\mathcal{A}$ semisederhana (atau terbelah), maka
    $\mathcal{A}$ disebut
    kategori abelian semisederhana (atau terbelah)\footnote{Definisi dalam
    literatur tidak seragam; sebagian penulis menyebut kategori abelian terbelah
    di sini sebagai kategori abelian semisederhana.}.
\end{definisi}

% segment-id: o014.li2.chapter2.simplicity-and-semisimplicity.p006
Banyak literatur mensyaratkan $I$ berhingga dalam definisi objek
semisederhana.

% segment-id: o014.li2.chapter2.simplicity-and-semisimplicity.lm003
\begin{lema}
    Jika $X=\bigoplus_{i=1}^nX_i$, dengan $X_1,\ldots,X_n$ objek sederhana,
maka $X$ berpanjang hingga, $\mathrm{JH}(X)=\left\{X_1,\ldots,X_n\right\}$
(dihitung dengan multiplisitas), dan $\ell(X)=n$.
\end{lema}
% segment-id: o014.li2.chapter2.simplicity-and-semisimplicity.pf004
\begin{bukti}
    Tinjau deret komposisi
$\bigoplus_{i=1}^nX_i\supsetneq\bigoplus_{i=1}^{n-1}X_i\supsetneq\cdots\supsetneq0$.
\end{bukti}

% segment-id: o014.li2.chapter2.simplicity-and-semisimplicity.rem001
\begin{catatan}
    Untuk objek semisederhana umum $X=\bigoplus_{i\in I}X_i$, sifat-sifat
Noetherian, Artinian, dan berpanjang hingga dari $X$ semuanya ekuivalen dengan
$I$ berhingga. Hal ini karena dengan menambahkan (atau menghapus) suku jumlah
langsung dari $I$, kita dapat dengan mudah membangun rantai naik ketat (atau
rantai menurun ketat) dalam $\mathrm{Sub}_X$.
\end{catatan}

% segment-id: o014.li2.chapter2.simplicity-and-semisimplicity.p007
Kita ingin memahami hubungan antara objek terbelah dan objek semisederhana.
Argumennya serupa dengan kasus modul \cite[Proposisi 6.11.4]{Li1}.

% segment-id: o014.li2.chapter2.simplicity-and-semisimplicity.pr001
\begin{proposisi}
    Misalkan $X$ merupakan objek terbelah. Maka subobjek dan objek kuosien dari
$X$ juga terbelah. Jika lebih lanjut $X$ merupakan objek Artinian, maka $X$
merupakan objek semisederhana berpanjang hingga.
\end{proposisi}
% segment-id: o014.li2.chapter2.simplicity-and-semisimplicity.pf005
\begin{bukti}
    Ambil barisan eksak pendek $0\to X'\to X\to X''\to0$. Kondisi tersebut
menyiratkan $X\simeq X'\oplus X''$, sehingga $X''$ dapat disisipkan sebagai
subobjek dari $X$. Jadi, cukup dibuktikan bahwa setiap
subobjek $X'$ dari $X$ terbelah. Diberikan $X'_0\subset X'$, terdapat
$Y\subset X$ sedemikian sehingga $X=X'_0\oplus Y$. Maka yang perlu
dibuktikan adalah
% segment-id: o014.li2.chapter2.simplicity-and-semisimplicity.q005
\[ X'=X'_0\oplus(Y\cap X'). \]
Ini merupakan penerapan Proposisi~\ref{prop:biproduct-sum-intersection}: di
satu sisi $X'_0\cap(Y\cap X')\subset X'_0\cap Y=0$; di sisi lain,
$\mathrm{Sub}_X$ merupakan kisi modular, sehingga
$X'_0+(Y\cap X')=X'\cap(Y+X'_0)=X'\cap X=X'$. Dengan demikian, persamaan di
atas terbukti.

    Selanjutnya, misalkan $X$ merupakan objek Artinian. Jika $X\neq0$, terdapat
subobjek tak nol minimal $X_1$; subobjek ini pasti sederhana, dan terdapat
dekomposisi jumlah langsung $X=X_1\oplus Y_1$. Perhatikan bahwa $Y_1$ tetap
merupakan objek Artinian yang terbelah; jika $Y_1\neq0$, lanjutkan dengan
menulis $Y_1=X_2\oplus Y_2$, dan seterusnya. Kondisi Artinian menjamin bahwa
rantai menurun ketat $X\supsetneq Y_1\supsetneq Y_2\cdots$ berhenti setelah
langkah berhingga, sehingga diperoleh dekomposisi yang diinginkan
$X=X_1\oplus\cdots\oplus X_n$.
\end{bukti}

% segment-id: o014.li2.chapter2.simplicity-and-semisimplicity.p008
Kita juga dapat menanyakan apakah kebalikannya berlaku: apakah objek
semisederhana terbelah.
Untuk kategori abelian umum, sekali lagi diperlukan asumsi keberhinggaan.

% segment-id: o014.li2.chapter2.simplicity-and-semisimplicity.pr002
\begin{proposisi}\label{prop:ss-decomp}
    Untuk suatu objek $X$, tinjau sifat-sifat berikut.
    % segment-id: o014.li2.chapter2.simplicity-and-semisimplicity.l002
    \begin{enumerate}[(i)]
        % segment-id: o014.li2.chapter2.simplicity-and-semisimplicity.i003
        \item $X=\sum_{Y\in\mathcal{F}}Y$, dengan $\mathcal{F}$ suatu himpunan
        bagian berhingga dari $\mathrm{Sub}_X$ dan setiap $Y\in\mathcal{F}$
        merupakan objek sederhana;
        % segment-id: o014.li2.chapter2.simplicity-and-semisimplicity.i004
        \item $X=\bigoplus_{Y\in\mathcal{F}}Y$, dengan $\mathcal{F}$ suatu
        himpunan bagian berhingga dari $\mathrm{Sub}_X$ dan setiap
        $Y\in\mathcal{F}$ merupakan objek sederhana;
        % segment-id: o014.li2.chapter2.simplicity-and-semisimplicity.i005
        \item $X$ terbelah.
    \end{enumerate}
    Berlaku (i) $\implies$ (ii) $\implies$ (iii).
\end{proposisi}
% segment-id: o014.li2.chapter2.simplicity-and-semisimplicity.pf006
\begin{bukti}
    Ulangi argumen untuk (i) $\implies$ (ii) $\implies$ (iii) dalam
\cite[Proposisi 6.11.4]{Li1}.
\end{bukti}

% segment-id: o014.li2.chapter2.simplicity-and-semisimplicity.rem002
\begin{catatan}\label{rem:Grothendieck-ss-decomp}
    Untuk memperluas pernyataan Proposisi~\ref{prop:ss-decomp} ke himpunan bagian
tak hingga $\mathcal{F}\subset\mathrm{Sub}_X$, kita harus mensyaratkan bahwa
$\mathcal{A}$ merupakan kategori Grothendieck yang akan diperkenalkan dalam
\S~\sourcecrossref{sec:Grothendieck-cat}{2.10}. Dengan demikian, setiap objek semisederhana
dalam kategori Grothendieck terbelah, dan kategori Grothendieck semisederhana
otomatis terbelah. Setelah definisi terkait dikuasai, argumennya serupa dengan
kasus modul, sehingga hal ini diserahkan sebagai latihan dalam
bab ini.
\end{catatan}
